% GENERATED FILE. Edit canonical lessons, then run npm run generate:books.
% canonical-source-hash: fnv1a64:fde1a098cc29527f
% canonical-lessons: TA-C547-kaikutu, TA-C547-neritu, TA-C547-utteci, TA-C547-niti, TA-C547-nal-kuri

\chapter{When Plans Work Out}
\label{ch:ta-when-plans-work-out}
\hlchaptermodality{547}

\begin{chapteropening}
\textbf{By the end of this chapter:} \emph{I can talk about plans working out and dates being fixed.}

Complete the last lesson of chapter 547: I can talk about plans working out and dates being fixed.
\end{chapteropening}

\section[\texorpdfstring{\ta{கைகூடு} (kaikūṭu)}{கைகூடு (kaikūṭu)}]{\ta{கைகூடு} (kaikūṭu) — to work out; to come true}

\label{lesson:TA-C547-kaikutu}



\begin{quote}
Before the new one: say the Tamil for a stumbling block, then the Tamil for to hold out.
\end{quote}

\subsection*{You'll want to know: \ta{கைகூடு}}
\textbf{\ta{கைகூடு}} — \emph{kaikūṭu} — \textquotedblleft{}to work out; to come true\textquotedblright{}.

\begin{quote}
\ta{எல்லாம்} \ta{சரியாக} \ta{நடந்தால்} \ta{எங்கள்} \ta{கனவு} \ta{கைகூடும்}.
\end{quote}

\begin{quote}
\emph{ellām cariyāka naṭantāl eṅkaḷ kaṉavu kaikūṭum.}
\end{quote}

\begin{quote}
If everything goes right, our dream will come true.
\end{quote}

Past (it …): \ta{கைகூடியது}, \emph{kaikūṭiyatu}.

The plan or wish is the subject, so the past is shown with \ta{அது}: \ta{முயற்சி} \ta{கைகூடியது}, the effort succeeded.

\begin{cousinweb}
\textbf{\ta{கைகூடு}} (\emph{kaikūṭu}) is a native compound: \ta{கை} (\emph{kai}), \textquotedblleft{}a hand\textquotedblright{}, and \ta{கூடு} (\emph{kūṭu}), \textquotedblleft{}to come together, to be achieved\textquotedblright{} (DEDR 1882): to come into one's hand.
\end{cousinweb}

\begin{grammarlens}[title={What you've built}]
One more word for saying what happens if.
\end{grammarlens}

\subsection*{Your turn}
\begin{itemize}
\raggedright
  \item \emph{Say it:} \emph{kaikūṭu}
  \item \emph{Say it:} \emph{kaikūṭu}, once more
  \item \emph{Say it:} \emph{ellām cariyāka naṭantāl eṅkaḷ kaṉavu kaikūṭum.} — If everything goes right, our dream will come true.
  \item \emph{Recall:} say the Tamil for a stumbling block, then the Tamil for to hold out, then say \emph{kaikūṭu} again
\end{itemize}

\begin{tcolorbox}[breakable,colback=teal!4,colframe=teal!35!black,title={Before you move on}]
What does \ta{கைகூடு} mean? (\textquotedblleft{}To work out; to come true\textquotedblright{}.) Say it once more.
\end{tcolorbox}
\section[\texorpdfstring{\ta{நேரிடு} (nēriṭu)}{நேரிடு (nēriṭu)}]{\ta{நேரிடு} (nēriṭu) — to befall; to happen}

\label{lesson:TA-C547-neritu}



\begin{quote}
Before the new one: say the Tamil for to hold out, then the Tamil for to work out.
\end{quote}

\subsection*{You'll want to know: \ta{நேரிடு}}
\textbf{\ta{நேரிடு}} — \emph{nēriṭu} — \textquotedblleft{}to befall; to happen\textquotedblright{}.

\begin{quote}
\ta{கவனமாக} \ta{இல்லாவிட்டால்} \ta{விபத்து} \ta{நேரிடலாம்}.
\end{quote}

\begin{quote}
\emph{kavaṉamāka illāviṭṭāl vipattu nēriṭalām.}
\end{quote}

\begin{quote}
If you are not careful, an accident may happen.
\end{quote}

Past (it …): \ta{நேரிட்டது}, \emph{nēriṭṭatu}.

Of bad things, with \ta{அது} as subject: \ta{விபத்து} \ta{நேரிட்டது}, an accident happened. After an infinitive it means \textquotedblleft{}to have to\textquotedblright{}: \ta{போக} \ta{நேரிட்டது}.

\begin{cousinweb}
\textbf{\ta{நேரிடு}} (\emph{nēriṭu}) is a native compound: \ta{நேர்} (\emph{nēr}), \textquotedblleft{}to meet, to come near\textquotedblright{} (DEDR 3770), and \ta{இடு}. It is not built on \ta{நேரம்}, \textquotedblleft{}time\textquotedblright{}.
\end{cousinweb}

\begin{grammarlens}[title={What you've built}]
One more word for saying what happens if.
\end{grammarlens}

\subsection*{Your turn}
\begin{itemize}
\raggedright
  \item \emph{Say it:} \emph{nēriṭu}
  \item \emph{Say it:} \emph{nēriṭu}, once more
  \item \emph{Say it:} \emph{kavaṉamāka illāviṭṭāl vipattu nēriṭalām.} — If you are not careful, an accident may happen.
  \item \emph{Recall:} say the Tamil for to hold out, then the Tamil for to work out, then say \emph{nēriṭu} again
\end{itemize}

\begin{tcolorbox}[breakable,colback=teal!4,colframe=teal!35!black,title={Before you move on}]
What does \ta{நேரிடு} mean? (\textquotedblleft{}To befall; to happen\textquotedblright{}.) Say it once more.
\end{tcolorbox}
\section[\texorpdfstring{\ta{உத்தேசி} (uttēci)}{உத்தேசி (uttēci)}]{\ta{உத்தேசி} (uttēci) — to intend; to estimate}

\label{lesson:TA-C547-utteci}



\begin{quote}
Before the new one: say the Tamil for to work out, then the Tamil for to befall.
\end{quote}

\subsection*{You'll want to know: \ta{உத்தேசி}}
\textbf{\ta{உத்தேசி}} — \emph{uttēci} — \textquotedblleft{}to intend; to estimate\textquotedblright{}.

\begin{quote}
\ta{அடுத்த} \ta{ஆண்டு} \ta{வீடு} \ta{கட்ட} \ta{உத்தேசித்திருக்கிறோம்}.
\end{quote}

\begin{quote}
\emph{aṭutta āṇṭu vīṭu kaṭṭa uttēcittirukkiṟōm.}
\end{quote}

\begin{quote}
We intend to build a house next year.
\end{quote}

Past (I …): \ta{உத்தேசித்தேன்}, \emph{uttēcittēṉ}.

The plan is an infinitive: \ta{கட்ட} \ta{உத்தேசி}, to intend to build.

\begin{cousinweb}
\textbf{\ta{உத்தேசி}} (\emph{uttēci}) is built on Sanskrit \emph{uddēśa}, \textquotedblleft{}pointing out, explaining\textquotedblright{} (CDIAL 1998), made a Tamil verb with -\ta{இ}.
\end{cousinweb}

\begin{grammarlens}[title={What you've built}]
One more word for saying what happens if.
\end{grammarlens}

\subsection*{Your turn}
\begin{itemize}
\raggedright
  \item \emph{Say it:} \emph{uttēci}
  \item \emph{Say it:} \emph{uttēci}, once more
  \item \emph{Say it:} \emph{aṭutta āṇṭu vīṭu kaṭṭa uttēcittirukkiṟōm.} — We intend to build a house next year.
  \item \emph{Recall:} say the Tamil for to work out, then the Tamil for to befall, then say \emph{uttēci} again
\end{itemize}

\begin{tcolorbox}[breakable,colback=teal!4,colframe=teal!35!black,title={Before you move on}]
What does \ta{உத்தேசி} mean? (\textquotedblleft{}To intend; to estimate\textquotedblright{}.) Say it once more.
\end{tcolorbox}
\section[\texorpdfstring{\ta{நீடி} (nīṭi)}{நீடி (nīṭi)}]{\ta{நீடி} (nīṭi) — to persist; to last}

\label{lesson:TA-C547-niti}



\begin{quote}
Before the new one: say the Tamil for to befall, then the Tamil for to intend.
\end{quote}

\subsection*{You'll want to know: \ta{நீடி}}
\textbf{\ta{நீடி}} — \emph{nīṭi} — \textquotedblleft{}to persist; to last\textquotedblright{}.

\begin{quote}
\ta{மழை} \ta{இன்னும்} \ta{இரண்டு} \ta{நாட்கள்} \ta{நீடிக்கும்}.
\end{quote}

\begin{quote}
\emph{maḻai iṉṉum iraṇṭu nāṭkaḷ nīṭikkum.}
\end{quote}

\begin{quote}
The rain will last two more days.
\end{quote}

Past (it …): \ta{நீடித்தது}, \emph{nīṭittatu}.

Usually with a thing or a time as subject, so the past is shown with \ta{அது}.

\begin{cousinweb}
\textbf{\ta{நீடி}} (\emph{nīṭi}) is native Dravidian: DEDR 3692 gives \ta{நீள்} (\emph{nīḷ}), \textquotedblleft{}to be long\textquotedblright{}, and \ta{நீடு}, \textquotedblleft{}to grow long, to be extended through time\textquotedblright{}.
\end{cousinweb}

\begin{grammarlens}[title={What you've built}]
One more word for saying what happens if.
\end{grammarlens}

\subsection*{Your turn}
\begin{itemize}
\raggedright
  \item \emph{Say it:} \emph{nīṭi}
  \item \emph{Say it:} \emph{nīṭi}, once more
  \item \emph{Say it:} \emph{maḻai iṉṉum iraṇṭu nāṭkaḷ nīṭikkum.} — The rain will last two more days.
  \item \emph{Recall:} say the Tamil for to befall, then the Tamil for to intend, then say \emph{nīṭi} again
\end{itemize}

\begin{tcolorbox}[breakable,colback=teal!4,colframe=teal!35!black,title={Before you move on}]
What does \ta{நீடி} mean? (\textquotedblleft{}To persist; to last\textquotedblright{}.) Say it once more.
\end{tcolorbox}
\section[\texorpdfstring{\ta{நாள்} \ta{குறி} (nāḷ kuṟi)}{நாள் குறி (nāḷ kuṟi)}]{\ta{நாள்} \ta{குறி} (nāḷ kuṟi) — to fix a date}

\label{lesson:TA-C547-nal-kuri}



\begin{quote}
Before the new one: say the Tamil for to intend, then the Tamil for to persist.
\end{quote}

\subsection*{You'll want to know: \ta{நாள்} \ta{குறி}}
\textbf{\ta{நாள்} \ta{குறி}} — \emph{nāḷ kuṟi} — \textquotedblleft{}to fix a date\textquotedblright{}.

\begin{quote}
\ta{இரு} \ta{குடும்பங்களும்} \ta{திருமணத்துக்கு} \ta{நாள்} \ta{குறித்தார்கள்}.
\end{quote}

\begin{quote}
\emph{iru kuṭumpaṅkaḷum tirumaṇattukku nāḷ kuṟittārkaḷ.}
\end{quote}

\begin{quote}
The two families fixed a date for the wedding.
\end{quote}

Past (I …): \ta{நாள்} \ta{குறித்தேன்}, \emph{nāḷ kuṟittēṉ}.

The event takes -\ta{க்கு}: \ta{திருமணத்துக்கு} \ta{நாள்} \ta{குறி}.

\begin{cousinweb}
\textbf{\ta{நாள்} \ta{குறி}} (\emph{nāḷ kuṟi}) joins the native \ta{நாள்} (\emph{nāḷ}), \textquotedblleft{}a day\textquotedblright{} (DEDR 3656), and \ta{குறி} (\emph{kuṟi}), \textquotedblleft{}to mark, to fix\textquotedblright{} (DEDR 1847).
\end{cousinweb}

\begin{grammarlens}[title={What you've built}]
One more word for saying what happens if.
\end{grammarlens}

\subsection*{Your turn}
\begin{itemize}
\raggedright
  \item \emph{Say it:} \emph{nāḷ kuṟi}
  \item \emph{Say it:} \emph{nāḷ kuṟi}, once more
  \item \emph{Say it:} \emph{iru kuṭumpaṅkaḷum tirumaṇattukku nāḷ kuṟittārkaḷ.} — The two families fixed a date for the wedding.
  \item \emph{Recall:} say the Tamil for to intend, then the Tamil for to persist, then say \emph{nāḷ kuṟi} again
\end{itemize}

\begin{tcolorbox}[breakable,colback=teal!4,colframe=teal!35!black,title={Before you move on}]
What does \ta{நாள்} \ta{குறி} mean? (\textquotedblleft{}To fix a date\textquotedblright{}.) Say it once more.
\end{tcolorbox}
